\partchapter{3}{第三部分\quad 完成情况及性能参数}

\Needspace{85pt}

\section{整体介绍}

当前 PYNQ release 已在领航者 ZYNQ7020 上运行：PYNQ 2.7.0、Python 3.8.2，PS7 FCLK0 配置 100 MHz。 Python 通过 AXI-Lite/UART mailbox 调用 PL 内 RV32，同一 overlay 可切换 SW/HW/ISE，无需重载 bitstream。板级容量内 KAT、Notebook、warm benchmark 与错误恢复均已留证；小时级稳定性与独立频率尚未测量。

\figpdf{3-1}{fig_3_1}{系统实物正面照片}{docs/tech\_report/figs/board\_photo.jpg， 可见领航者 ATK-DF7010/7020P 丝印与上电指示。照片不证明拍摄时运行的 bitstream、芯片速度等级或长时稳定性； 背面、核心芯片完整丝印、跳线与 SD 镜像记录仍需补齐。}

{\small
\renewcommand{\thetable}{3-1}
\begin{longtable}{@{}>{\raggedright\arraybackslash}p{\dimexpr 0.3265\linewidth-2\tabcolsep\relax}>{\raggedright\arraybackslash}p{\dimexpr 0.6735\linewidth-2\tabcolsep\relax}@{}}
\caption{整体完成情况}\label{tab:3-1}\\
\toprule
\textbf{项目} & \textbf{已有证据} \\
\midrule
\endfirsthead
\multicolumn{2}{@{}l}{\footnotesize 续表 3-1}\\
\toprule
\textbf{项目} & \textbf{已有证据} \\
\midrule
\endhead
当前 KAT & 三模式各 289/289，坏 Tag 各 289/289 \\
当前加密中位周期 & SW/HW/ISE=23342/1439/3995 \\
错误恢复 & 9/9 checks、7/7 reset；每模式 50 次加解密往返 \\
实现时序 & WNS +0.034 ns、TNS 0，DRC 0 Error/0 Critical \\
当前 bit/hwh & 4045688 B / 148833 B，配套身份与哈希 \\
\bottomrule
\end{longtable}
}

\Needspace{85pt}

\section{工程成果}

当前成果分为硬件/固件源、配套 release 与原始验证数据。工程入口为 board/vivado\_zynq7000\_xc7z020\_pynq/build.tcl；BRAM COE 由服务固件 build.sh 生成。实际使用的 top、part、源集、IP 与源文件哈希见发布 manifest，不能以不存在的脚本或目录代替。

{\small
\renewcommand{\thetable}{3-2}
\begin{longtable}{@{}>{\raggedright\arraybackslash}p{\dimexpr 0.1960\linewidth-2\tabcolsep\relax}>{\raggedright\arraybackslash}p{\dimexpr 0.3753\linewidth-2\tabcolsep\relax}>{\raggedright\arraybackslash}p{\dimexpr 0.4287\linewidth-2\tabcolsep\relax}@{}}
\caption{工程成果清单}\label{tab:3-2}\\
\toprule
\textbf{产物} & \textbf{实际位置} & \textbf{版本与范围} \\
\midrule
\endfirsthead
\multicolumn{3}{@{}l}{\footnotesize 续表 3-2}\\
\toprule
\textbf{产物} & \textbf{实际位置} & \textbf{版本与范围} \\
\midrule
\endhead
当前 overlay/固件 & board/vivado\_\allowbreak{}zynq7000\_\allowbreak{}xc7z020\_\allowbreak{}pynq/release/ & bit/\allowbreak{}hwh/\allowbreak{}bin/\allowbreak{}coe/\allowbreak{}elf/\allowbreak{}manifest/\allowbreak{}sha256sums \\
当前实现与回归 & build/\allowbreak{}reports/\allowbreak{}pynq\_\allowbreak{}final\_\allowbreak{}20261028/ & timing.rpt、utilization.rpt、drc.rpt、summary.json \\
当前容量内 KAT & data/pynq/board\_\allowbreak{}20261028/ & 官方 289 组 $\times$ 3 模式，含错误 Tag 拒绝 \\
当前固定输入采样 & data/\allowbreak{}pynq/\allowbreak{}notebook\_\allowbreak{}20261029\_\allowbreak{}130920/ & 各档 30 样本与 PS 墙钟 \\
当前错误恢复 & data/pynq/error\_\allowbreak{}gate\_\allowbreak{}20261029\_\allowbreak{}134023/ & JSON、log、identity、已执行 ipynb \\
历史 PL-only 镜像 & build/\allowbreak{}local/\allowbreak{}board\_\allowbreak{}bitstream/ & 三镜像各 4045687 B；本地备份不是 clone 保证 \\
历史 PL-only 日志 & data/\allowbreak{}ascon/\allowbreak{}results/\allowbreak{}board/\allowbreak{} & uart\_\allowbreak{}20261026\_\allowbreak{}sw\_\allowbreak{}and\_\allowbreak{}ise.txt、uart\_\allowbreak{}20261027\_\allowbreak{}mmio.txt \\
\bottomrule
\end{longtable}
}

clean-clone 与最终证据包仍需独立验收；现有工件校验通过不等于源到板全流程已干净重建。

\Needspace{85pt}

\subsection{资源占用}

当前 \texttt{pynq\_\allowbreak{}wrapper} 的完整资源原表为 \texttt{build/\allowbreak{}reports/\allowbreak{}pynq\_\allowbreak{}final\_\allowbreak{}20261028/\allowbreak{}utilization.\allowbreak{}rpt}，数字由证据脚本抽取并列在附录 C，不用历史 PL-only 数字代替。

{\small
\renewcommand{\thetable}{3-3}
\begin{longtable}{@{}>{\raggedright\arraybackslash}p{\dimexpr 0.2583\linewidth-2\tabcolsep\relax}>{\raggedright\arraybackslash}p{\dimexpr 0.3663\linewidth-2\tabcolsep\relax}>{\raggedright\arraybackslash}p{\dimexpr 0.3754\linewidth-2\tabcolsep\relax}@{}}
\caption{当前 PYNQ 实现时序与 DRC}\label{tab:3-3}\\
\toprule
\textbf{当前 PYNQ 指标} & \textbf{结果} & \textbf{原始报告} \\
\midrule
\endfirsthead
\multicolumn{3}{@{}l}{\footnotesize 续表 3-3}\\
\toprule
\textbf{当前 PYNQ 指标} & \textbf{结果} & \textbf{原始报告} \\
\midrule
\endhead
WNS / TNS & +0.034 ns / 0 & timing.rpt \\
WHS / THS & +0.020 ns / 0 & timing.rpt \\
setup / hold 失败 & 均 0/62861 & timing.rpt \\
内部未约束端点 & 0 & timing.rpt 的 check\_\allowbreak{}timing 段 \\
DRC & 0 Error / 0 Critical Warning & drc.rpt \\
\bottomrule
\end{longtable}
}

整机资源为 LUT 14269、FF 19691、BRAM 34.5 tile、DSP 4。 DRC 默认报告有 53 项 Warning，其中包含规则条数上限提示；不能将该显示数当作所有潜在 Warning 的总数，更不能写成``DRC 无任何警告''。DSP 流水、RAM 输出寄存器等提示保留在原始报告。

历史 PL-only 整机资源为 LUT 12426、FF 16156、BRAM 34.5 tile、DSP 4，WHS +0.021 ns。原件位于 \texttt{build/\allowbreak{}local/\allowbreak{}proj\_\allowbreak{}backup\_\allowbreak{}first\_\allowbreak{}time\_\allowbreak{}20261027/\allowbreak{}first\_\allowbreak{}time.\allowbreak{}runs/\allowbreak{}impl\_\allowbreak{}1/\allowbreak{}}，不是当前工程报告。这些是整机总资源，不据此推导每条加速路径的独立面积或功耗。

\figpdf{3-2}{fig_3_6}{Vivado 实现后器件视图}{打开 Implemented Design 的 Device 窗口， 展示 PS7 与 PL 内 CPU、Cache、Ascon 等逻辑在 XC7Z020 上的布局。}

\figpdf{3-3}{fig_3_7}{Vivado 资源利用率报告}{Report Utilization 的汇总页， 应可见 LUT 14269、FF 19691、BRAM 34.5、DSP 4。}

\Needspace{85pt}

\subsection{可复现性与产物冻结}

{\small
\renewcommand{\thetable}{3-4}
\begin{longtable}{@{}>{\raggedright\arraybackslash}p{\dimexpr 0.3241\linewidth-2\tabcolsep\relax}>{\raggedright\arraybackslash}p{\dimexpr 0.6759\linewidth-2\tabcolsep\relax}@{}}
\caption{发布身份记录}\label{tab:3-4}\\
\toprule
\textbf{对象} & \textbf{当前记录} \\
\midrule
\endfirsthead
\multicolumn{2}{@{}l}{\footnotesize 续表 3-4}\\
\toprule
\textbf{对象} & \textbf{当前记录} \\
\midrule
\endhead
release source commit & 86dc10f13e96adbb360f1adfe0c4891ce2f8e2a7；source\_\allowbreak{}dirty=true \\
板端 benchmark Git & ac9143d5c6acc9e7ac61f907a7e802eb6b35cddd \\
pynq.bit & 4045688 B；SHA-256 前缀 6e7f9320 \\
pynq.hwh & 148833 B；SHA-256 前缀 3ebd2640 \\
固件 & pynq-svc-v1；ELF/BIN/COE 逐项哈希见 manifest \\
身份核验 & 板端 release\_\allowbreak{}hashes.log 与 identity.json；bit/hwh/固件必须配套 \\
\bottomrule
\end{longtable}
}

bit/hwh 来自同一 BD 实现；BRAM MIF 与最终 BIN 的 128 KiB 补零镜像有逐字节一致性证据。 manifest 的 source\_sha256 是重建输入核对依据；source\_dirty=true 表示只给 Git commit 不足以冻结全部输入。 manifest 中 hardware\_validation=NOT\_RUN 是构建时快照，后续板测应引用日期明确的外部证据，本次不修改已发布 manifest 来追写状态。

从干净环境重建并运行的 D05 仍未完成；本文提供实际入口与现有记录，不承诺 clean clone 一键全部验收。复建不要求 bitstream 字节必然相同，应分别核对源输入、功能、时序、资源和新产物身份。

\evid{证据 board/\allowbreak{}vivado\_\allowbreak{}zynq7000\_\allowbreak{}xc7z020\_\allowbreak{}pynq/\allowbreak{}release/\allowbreak{}manifest.\allowbreak{}json；board/\allowbreak{}vivado\_\allowbreak{}zynq7000\_\allowbreak{}xc7z020\_\allowbreak{}pynq/\allowbreak{}validation\_\allowbreak{}20261028.\allowbreak{}md；validation\_\allowbreak{}20261029.\allowbreak{}md；build/\allowbreak{}reports/\allowbreak{}pynq\_\allowbreak{}final\_\allowbreak{}20261028/\allowbreak{}；docs/\allowbreak{}timing\_\allowbreak{}acb\_\allowbreak{}validation.\allowbreak{}md}

\Needspace{85pt}

\section{特性成果}

本节按''时序 $\rightarrow$ IO $\rightarrow$ 功能 $\rightarrow$ 性能 $\rightarrow$ 完成度''五项列出已达成的量化成果。前三项回答''这套设计能不能稳定工作''，后两项回答''它有多快、还有哪些没做''。板上测量、仿真与实现报告分别标注；未测量项记为 \texttt{未测量} 或 \texttt{未验证}，不以估计值代替。

\Needspace{85pt}

\subsection{时序收敛结果}

本节为带原始记录的研制过程，不将历史余量称为当前 release 余量。历史默认策略 WNS $-$0.942 ns，33 组策略中最优 Performance\_ExploreWithRemap 为 $-$0.118 ns、TNS $-$1.008 ns、22 个失败端点。记录见 \texttt{data/\allowbreak{}ascon/\allowbreak{}results/\allowbreak{}strategy\_\allowbreak{}sweep\_\allowbreak{}20261017/\allowbreak{}strategy\_\allowbreak{}sweep\_\allowbreak{}summary.\allowbreak{}json}。它说明已测策略的效果，不证明所有可能策略的理论上限。

对 fix\_r1 快照的端点归因识别出 ISE 私有状态复位扇出、乘法器组合路径及 D-cache 请求通路三类问题。对应结构改造为：局部复位注册化、乘法器输入输出注册与 ack/cancel 握手、MISS\_CAPTURE 与 store/MMIO 绕过。归因分类不是互斥端点总数之和，不以分类数相加覆盖所有端点。

{\small
\renewcommand{\thetable}{3-5}
\begin{longtable}{@{}>{\raggedright\arraybackslash}p{\dimexpr 0.4531\linewidth-2\tabcolsep\relax}>{\raggedright\arraybackslash}p{\dimexpr 0.1435\linewidth-2\tabcolsep\relax}>{\raggedright\arraybackslash}p{\dimexpr 0.1435\linewidth-2\tabcolsep\relax}>{\raggedright\arraybackslash}p{\dimexpr 0.2598\linewidth-2\tabcolsep\relax}@{}}
\caption{时序收敛历史记录}\label{tab:3-5}\\
\toprule
\textbf{历史 run} & \textbf{WNS ns} & \textbf{TNS ns} & \textbf{setup 失败端点} \\
\midrule
\endfirsthead
\multicolumn{4}{@{}l}{\footnotesize 续表 3-5}\\
\toprule
\textbf{历史 run} & \textbf{WNS ns} & \textbf{TNS ns} & \textbf{setup 失败端点} \\
\midrule
\endhead
A+C+B / ExplorePostRoutePhysOpt & +0.001 & 0 & 0 \\
A+C+B / ExploreWithRemap & +0.007 & 0 & 0 \\
PL-only 三镜像报告 & +0.087 & 0 & 0/55572 \\
当前 PYNQ release & +0.034 & 0 & 0/62861 \\
\bottomrule
\end{longtable}
}

前两项为历史策略快照，第三项为历史板测镜像实现，第四项为独立 PYNQ BD 实现；不能将相邻行之差全归因于某一补丁。当前余量仅 34 ps，任何功能、约束或 IP 变化都须重新验证。

\figpdf{3-4}{fig_3_2}{历史 PL-only 时序收敛记录}{默认策略、已测策略最优、 结构改造与历史板测镜像的 WNS；不替代当前 PYNQ 实现报告。}

\Needspace{85pt}

\subsection{时钟与 IO 的适用范围}\label{ux65f6ux949fux4e0e-io-ux7684ux9002ux7528ux8303ux56f4}

当前 PYNQ 通路位于 PL 内部，通过 PS 的 AXI-Lite 窗口交互，没有沿用历史板外 UART、LED、按键 IO XDC。当前报告内部未约束端点为 0，未列出 no\_clock/no\_input\_delay/no\_output\_delay 问题。这不等于测得了板外 IO 延迟，也不能把别的源集的 IO 覆盖检查作为本工程签核。

历史板外 UART 的约束分析仅作为工程资料保留，不列为当前 release 的新增成果。 FCLK0 配置为 100 MHz，但独立频率测量仍为 not\_measured；本报告不据配置值声称实测 Fmax。时序报告中的 WNS、TNS、WHS、THS 与路径覆盖须联合检查，不能只展示一个正 WNS。

\Needspace{85pt}

\subsection{功能验证}

\Needspace{85pt}

\paragraph{3.3.3.1 验证体系与层级划分}

{\small
\renewcommand{\thetable}{3-6}
\begin{longtable}{@{}>{\raggedright\arraybackslash}p{\dimexpr 0.0984\linewidth-2\tabcolsep\relax}>{\raggedright\arraybackslash}p{\dimexpr 0.1889\linewidth-2\tabcolsep\relax}>{\raggedright\arraybackslash}p{\dimexpr 0.2384\linewidth-2\tabcolsep\relax}>{\raggedright\arraybackslash}p{\dimexpr 0.2735\linewidth-2\tabcolsep\relax}>{\raggedright\arraybackslash}p{\dimexpr 0.2008\linewidth-2\tabcolsep\relax}@{}}
\caption{验证体系与层级划分}\label{tab:3-6}\\
\toprule
\textbf{层级} & \textbf{对象} & \textbf{手段} & \textbf{判据} & \textbf{状态} \\
\midrule
\endfirsthead
\multicolumn{5}{@{}l}{\footnotesize 续表 3-6}\\
\toprule
\textbf{层级} & \textbf{对象} & \textbf{手段} & \textbf{判据} & \textbf{状态} \\
\midrule
\endhead
L0 标准一致性 & C 与 Python 参考实现 & 独立下载的官方 KAT 逐字节比对 & 1,089 组加密、解密与负向用例全部通过 & 完成 \\
L1 模块级 RTL & 置换器、AEAD 适配器、Cache、ISE & Icarus/XSim 定向 TB 与 golden 状态比对 & 与锁定参考逐字节一致，事务契约成立 & 完成 \\
L2 系统级 & 整个 SoC 与真实固件 & \texttt{run\_\allowbreak{}\allowbreak{}tests.\allowbreak{}py} 自动回归（smoke 与 regression） & 报告 JSON/Markdown，非零退出即判失败 & 完成（含 NOT\_\allowbreak{}RUN 项） \\
L3 板级 & 三份镜像与真实芯片 & 串口 KAT 与 benchmark 采集 & SW 6 组、ISE 与 MMIO 各 4 组官方向量，含篡改拒绝与周期数 & 完成 \\
L4 PS/PL 协同 & 当前 release、服务固件与 host API & 容量内 KAT、Notebook、warm benchmark、错误恢复 & 三模式各 289/289、9/9 错误门禁；simulated=false & 有限范围通过（见 §2.\allowbreak{}3.\allowbreak{}3.\allowbreak{}6-\allowbreak{}-\allowbreak{}2.\allowbreak{}3.\allowbreak{}3.\allowbreak{}11） \\
\bottomrule
\end{longtable}
}

\figpdf{3-5}{fig_3_4}{五层验证体系与证据流}{host $\rightarrow$ RTL $\rightarrow$ SoC 仿真 $\rightarrow$ 历史 PL-only 板测 $\rightarrow$ 当前 PYNQ KAT、benchmark 与错误恢复，各层单列范围。}

\Needspace{85pt}

\paragraph{3.3.3.2 L0：标准一致性与向量冻结}\label{l0ux6807ux51c6ux4e00ux81f4ux6027ux4e0eux5411ux91cfux51bbux7ed3}

该层的作用是将标准固化为不可变的字节序列，作为 RTL 与板级验证的唯一参照。

\begin{itemize}
\tightlist
\item
  \textbf{标准版本}：NIST SP 800-232 Final（2025-08-13 发布），本地 PDF 哈希 \texttt{1be9bb9a…}。旧版 Ascon v1.2 的 Ascon-128 与 Ascon-128a 在端序、IV 与 padding 上均不同，不可沿用其 KAT；
\item
  \textbf{参考实现}：ascon-c 固定 commit \texttt{446347f2}（上游 \texttt{crypto\_\allowbreak{}aead/\allowbreak{}asconaead128/\allowbreak{}ref/\allowbreak{}}）与 pyascon 固定 commit \texttt{ed24e54a}。pyascon 的 README 仍标注 draft，故''是否对应 Final''不以文档表述判定，而以 ``通过完整官方 KAT''为判据；
\item
  \textbf{官方向量}：\texttt{LWC\_\allowbreak{}AEAD\_\allowbreak{}KAT\_\allowbreak{}128\_\allowbreak{}128.\allowbreak{}txt}（同 commit），共 1,089 组，即 PT 与 AD 长度在 0 至 32 的笛卡尔积，哈希 \texttt{bbbc3469…}。key、nonce、PT、AD 与 expected CT 均取自原始文件，未使用本仓库生成器自证；
\item
  \textbf{扩展向量}：\texttt{aead\_\allowbreak{}cases.\allowbreak{}json} 共 64 组，其中 54 组取自上游矩阵，10 组为锁定 C 参考的扩展用例（含空串、\texttt{a5} 与 \texttt{a500}、尾零、33/64/129-byte PT 及最长 81-byte AD），每组附带首、尾 CT 与 tag 翻转输入；
\item
  \textbf{置换向量}：22 个输入状态 $\times$（12 个单轮、P8、P12），共 308 组，源自锁定 C helper，并经 Python 二次校核；
\item
  \textbf{冻结机制}：\texttt{freeze\_\allowbreak{}vectors.\allowbreak{}py} 先令两个锁定参考通过完整独立下载的 KAT，再由锁定 C 参考产生扩展用例、由 Python 交叉检查；仅写入新建空目录，不覆写上游 KAT；两次重生成的三个数据文件字节完全一致。
\end{itemize}

实测结果如下表。

{\small
\renewcommand{\thetable}{3-7}
\begin{longtable}{@{}>{\raggedright\arraybackslash}p{\dimexpr 0.2952\linewidth-2\tabcolsep\relax}>{\raggedright\arraybackslash}p{\dimexpr 0.3495\linewidth-2\tabcolsep\relax}>{\raggedright\arraybackslash}p{\dimexpr 0.3553\linewidth-2\tabcolsep\relax}@{}}
\caption{L0 标准一致性检查}\label{tab:3-7}\\
\toprule
\textbf{检查项} & \textbf{结果} & \textbf{证据} \\
\midrule
\endfirsthead
\multicolumn{3}{@{}l}{\footnotesize 续表 3-7}\\
\toprule
\textbf{检查项} & \textbf{结果} & \textbf{证据} \\
\midrule
\endhead
两个锁定参考的完整 KAT 门禁 & \textbf{12,936 PASS / 0 FAIL}，exit 0 & \texttt{reference\_\allowbreak{}\allowbreak{}gate.\allowbreak{}json} \\
四个 backend 全套（含完整性检查） & \textbf{28,661 PASS / 0 FAIL}，exit 0 & \texttt{host\_\allowbreak{}\allowbreak{}validation.\allowbreak{}json} \\
harness 与 golden 接口 & 5 tests PASS，exit 0 & \texttt{harness.\allowbreak{}log} \\
生产 C raw 失败解密契约 & 23965 PASS / 0 FAIL & \texttt{failure\_\allowbreak{}\allowbreak{}contract.\allowbreak{}json} \\
修复前行为复现 & 失败输出保留为负向证据，不代表当前生产结果 & \texttt{failure\_\allowbreak{}\allowbreak{}contract\_\allowbreak{}\allowbreak{}original\_\allowbreak{}\allowbreak{}reproduction.\allowbreak{}json} \\
冻结数据重生成 & 4 文件 byte-identical，exit 0 & \texttt{audit\_\allowbreak{}\allowbreak{}snapshot.\allowbreak{}json} \\
\bottomrule
\end{longtable}
}

生产 C 实现已修复 bad tag 时的候选明文擦除与长度归零，锁定参考保持原字节。修复前失败与当前通过分别保存，不能以早期失败覆盖当前结果，也不能因 KAT 通过推断 raw buffer 已擦除。

\Needspace{85pt}

\paragraph{3.3.3.3 L1：模块级 RTL 验证}\label{l1ux6a21ux5757ux7ea7-rtl-ux9a8cux8bc1}

{\small
\renewcommand{\thetable}{3-8}
\begin{longtable}{@{}>{\raggedright\arraybackslash}p{\dimexpr 0.2450\linewidth-2\tabcolsep\relax}>{\raggedright\arraybackslash}p{\dimexpr 0.3343\linewidth-2\tabcolsep\relax}>{\raggedright\arraybackslash}p{\dimexpr 0.3517\linewidth-2\tabcolsep\relax}>{\raggedright\arraybackslash}p{\dimexpr 0.0690\linewidth-2\tabcolsep\relax}@{}}
\caption{L1 模块级 RTL 验证}\label{tab:3-8}\\
\toprule
\textbf{被测对象} & \textbf{TB} & \textbf{覆盖点} & \textbf{结果} \\
\midrule
\endfirsthead
\multicolumn{4}{@{}l}{\footnotesize 续表 3-8}\\
\toprule
\textbf{被测对象} & \textbf{TB} & \textbf{覆盖点} & \textbf{结果} \\
\midrule
\endhead
\texttt{ascon\_\allowbreak{}\allowbreak{}permutation}（单轮、P8、P12） & \texttt{tb\_\allowbreak{}\allowbreak{}ascon\_\allowbreak{}\allowbreak{}permutation}（golden memory） & 22 状态 $\times$ 14 操作 = 308 组，逐字节比对 & PASS \\
\texttt{ascon\_\allowbreak{}\allowbreak{}aead\_\allowbreak{}\allowbreak{}core}、\texttt{ascon\_\allowbreak{}\allowbreak{}aead\_\allowbreak{}\allowbreak{}2} & \texttt{tb\_\allowbreak{}\allowbreak{}ascon\_\allowbreak{}\allowbreak{}aead\_\allowbreak{}\allowbreak{}regression}、\texttt{tb\_\allowbreak{}\allowbreak{}ascon\_\allowbreak{}\allowbreak{}decrypt\_\allowbreak{}\allowbreak{}regression} & 全源加密与解密、篡改、长度边界 & PASS \\
AEAD ISE（\texttt{ascon\_\allowbreak{}\allowbreak{}ise}） & 定向 TB 与 EX 级集成 TB & 四指令编码、非法索引、取消、跨完成 hold & PASS \\
D-cache & 修订版 TB（\texttt{dcache\_\allowbreak{}\allowbreak{}tb\_\allowbreak{}\allowbreak{}revision.\allowbreak{}md}） & 命中、miss、填充、\texttt{MISS\_\allowbreak{}\allowbreak{}CAPTURE} 单拍、store 绕过、两路替换 & PASS \\
CPU 中断与异常 & 6 个真实 CPU 与 CLINT 中断场景 & 抢占、\texttt{mret} 恢复、\texttt{mepc} 正确性 & PASS \\
\bottomrule
\end{longtable}
}

关键做法为：\textbf{golden 数据不由 RTL 自身产生}。置换与 AEAD 的期望值均取自锁定 C 参考（编译为 shared library，由 Python 侧 ctypes 调用），RTL 仅提供''输入状态至输出状态''的映射，两侧在 Python 中比对。独立参考降低同源自证风险，但通过仅覆盖实际执行的向量与断言，不证明所有可能错误必被检出。

\figpdf{3-6}{fig_3_8}{Vivado 仿真波形：ISE 置换执行}{在 XSim 中打开 ISE 集成测试， 截取 custom-0 指令进入 EX、`ex\_stall\_req` 保持、置换完成后 `done` 与流水推进的一段波形。}

\Needspace{85pt}

\paragraph{3.3.3.4 L2：系统级自动回归}\label{l2ux7cfbux7edfux7ea7ux81eaux52a8ux56deux5f52}

入口为 \texttt{sim/\allowbreak{}board/\allowbreak{}run\_\allowbreak{}tests.\allowbreak{}py}，共两档：

{\small
\renewcommand{\thetable}{3-9}
\begin{longtable}{@{}>{\raggedright\arraybackslash}p{\dimexpr 0.2124\linewidth-2\tabcolsep\relax}>{\raggedright\arraybackslash}p{\dimexpr 0.3714\linewidth-2\tabcolsep\relax}>{\raggedright\arraybackslash}p{\dimexpr 0.3177\linewidth-2\tabcolsep\relax}>{\raggedright\arraybackslash}p{\dimexpr 0.0984\linewidth-2\tabcolsep\relax}@{}}
\caption{L2 系统级自动回归}\label{tab:3-9}\\
\toprule
\textbf{档位} & \textbf{内容} & \textbf{实测结果} & \textbf{耗时} \\
\midrule
\endfirsthead
\multicolumn{4}{@{}l}{\footnotesize 续表 3-9}\\
\toprule
\textbf{档位} & \textbf{内容} & \textbf{实测结果} & \textbf{耗时} \\
\midrule
\endhead
\texttt{-\allowbreak{}-\allowbreak{}mode smoke} & 快速入口：编译、镜像一致性、基础功能 & \textbf{6 PASS / 0 FAIL / 0 BLOCKED / 9 NOT\_\allowbreak{}RUN}，exit 0 & 18.739 s \\
\texttt{-\allowbreak{}-\allowbreak{}mode regression -\allowbreak{}-\allowbreak{}pynq-\allowbreak{}p0} & 指定 PYNQ P0 源集：三固件、host KAT、Mailbox、镜像、Cache 与定向中断 & \textbf{15 PASS / 0 FAIL / 0 BLOCKED / 1 NOT\_\allowbreak{}RUN}，exit 0 & 400.168 s \\
\bottomrule
\end{longtable}
}

NOT\_RUN 为显式结果，而非遗漏执行项：其对应''需真实板卡或需 PS7''的用例，在仿真环境中无法通过，故既不计 PASS 亦不计 FAIL。该三态设计的目的在于防止''不可执行的用例被标记为通过''。

本表只引用 2026-10-28 的归档 \texttt{smoke/\allowbreak{}summary.\allowbreak{}json} 与 \texttt{regression\_\allowbreak{}final/\allowbreak{}summary.\allowbreak{}json}， source\_head=0f5d49c，含 dirty 状态，不称为 10-29 当前 HEAD 的重新回归。产物含工具版本、命令、退出码与日志路径；历史 FAIL 与负向日志全部保留，不作删除。本档运行于 Windows 侧（Icarus/vvp 12 devel、GCC 13.2、 Python 3.14.3），与 host KAT 所用的 WSL 环境（Python 3.12.3、GCC 13.3.0，见 §3.3.3.2）分属两套工具链，故文中出现的多个 Python 版本各有其归属。

\Needspace{85pt}

\paragraph{3.3.3.5 L3：纯 PL 镜像的板级验证}\label{l3ux7eaf-pl-ux955cux50cfux7684ux677fux7ea7ux9a8cux8bc1}

\textbf{KAT 与篡改拒绝}（三份镜像，同一板卡、同一串口、同日执行）：

{\small
\renewcommand{\thetable}{3-10}
\begin{longtable}{@{}>{\raggedright\arraybackslash}p{\dimexpr 0.0804\linewidth-2\tabcolsep\relax}>{\raggedright\arraybackslash}p{\dimexpr 0.4908\linewidth-2\tabcolsep\relax}>{\raggedright\arraybackslash}p{\dimexpr 0.1306\linewidth-2\tabcolsep\relax}>{\raggedright\arraybackslash}p{\dimexpr 0.2981\linewidth-2\tabcolsep\relax}@{}}
\caption{L3 历史 PL-only 镜像 KAT 与篡改拒绝}\label{tab:3-10}\\
\toprule
\textbf{镜像} & \textbf{KAT 覆盖} & \textbf{篡改拒绝} & \textbf{附项} \\
\midrule
\endfirsthead
\multicolumn{4}{@{}l}{\footnotesize 续表 3-10}\\
\toprule
\textbf{镜像} & \textbf{KAT 覆盖} & \textbf{篡改拒绝} & \textbf{附项} \\
\midrule
\endhead
软件 & empty、1 B、16 B、16 B+AD、17 B、32 B（6 组官方向量） & 通过 & --- \\
ISE & empty、1 B、16 B、16 B+AD，另含解密往返 & 通过 & 解密往返一致性 \\
MMIO & 同 ISE 四组 & 通过 & \textbf{长度拒绝}（>16 B 显式拒绝） \\
\bottomrule
\end{longtable}
}

\textbf{benchmark}（16 字节加密、warm、\texttt{mcycle} 差值）：软件 20806 周期，ISE 2570 周期，MMIO 端到端 688 周期（核本体 103 周期）。同步采集的仿真值为 17353、2496、670 周期（核 99 周期），板级相对仿真的开销分别为 \textbf{+19.9\%、+3.0\%、+2.7\%}（核本体 +4.0\%），开销与访存密度正相关。

板级与仿真使用不同存储环境，不能仅凭输出差值认定某一等待路径就是全部成因。这些差异与存储模型、固件执行路径有关，但缺少分项实验，本文只保留观测，不声称已完成因果量化。

\textbf{原始证据}：串口日志按 XCOM 会话原始字节（\texttt{.\allowbreak{}XDat}）与去除非打印字符的可读版本（\texttt{.\allowbreak{}txt}）双份归档，并记录 SHA-256。MMIO 版日志开头存在约 2 行乱码（打开串口瞬间的帧错），banner 之后数据完整，两份可相互核对。覆盖组数以上表为准，不能把解密往返与篡改拒绝混记为额外官方向量。

\figpdf{3-7}{fig_3_5}{历史 PL-only 串口日志摘录}{SW/ISE 日志的 KAT 与 benchmark， 为文本重排而非界面截屏；正面实物照片另见系统实物图，其他现场素材待补。}

\Needspace{85pt}

\paragraph{3.3.3.6 缺陷发现与处置}

本节记录验证过程中发现的问题及其处置，按''发现、定性、修复、复验''组织，其信息量高于''通过项''列举。

{\small
\renewcommand{\thetable}{3-11}
\begin{longtable}{@{}>{\raggedright\arraybackslash}p{\dimexpr 0.0685\linewidth-2\tabcolsep\relax}>{\raggedright\arraybackslash}p{\dimexpr 0.3176\linewidth-2\tabcolsep\relax}>{\raggedright\arraybackslash}p{\dimexpr 0.1816\linewidth-2\tabcolsep\relax}>{\raggedright\arraybackslash}p{\dimexpr 0.3176\linewidth-2\tabcolsep\relax}>{\raggedright\arraybackslash}p{\dimexpr 0.1147\linewidth-2\tabcolsep\relax}@{}}
\caption{缺陷发现与处置}\label{tab:3-11}\\
\toprule
\textbf{序号} & \textbf{发现} & \textbf{定性} & \textbf{处理} & \textbf{状态} \\
\midrule
\endfirsthead
\multicolumn{5}{@{}l}{\footnotesize 续表 3-11}\\
\toprule
\textbf{序号} & \textbf{发现} & \textbf{定性} & \textbf{处理} & \textbf{状态} \\
\midrule
\endhead
1 & 标准审计所得 7 项静态缺陷（R01--R07）：长度由最高非零 bit 反推、\texttt{p12/\allowbreak{}p8} 未接 \texttt{start}、tag 输出被覆盖、多块指针推进错误、AD=8 时 pad 走错分支、decrypt 为空壳、busy 时 \texttt{start} 覆盖状态 & 算法与 RTL 接口不达标 & 改用显式长度，引入请求、接收与完成握手，独立 tag 输出与 block 索引；重写 \texttt{decrypt} 并增加 tag 校验 & 已修 \\
2 & 原生产 C 在 bad tag 后保留候选明文 & 失败 API 契约不满足项目要求 & 修生产 C 的擦除与长度归零；锁定参考不改，保留负向证据 & 23965 检查通过 \\
3 & D-cache miss 至总线请求为长组合路径（11 端点） & 时序 & 新增 \texttt{MISS\_\allowbreak{}\allowbreak{}CAPTURE}；store 与 MMIO 绕过 & 已修 \\
4 & 乘法器组合深度（48 端点、$-$0.105 ns） & 时序 & 输入与输出寄存器化，引入 \texttt{mul\_\allowbreak{}\allowbreak{}ack} 与 \texttt{cancel} & 已修 \\
5 & ISE 私有状态复位扇出（93 端点、$-$0.128 ns） & 时序 & 局部复位寄存器化，取消当拍清架构态 & 已修 \\
6 & 历史仿真与板级周期差达 +19.9\% & 观测，分项归因未完成 & 各自保留原始值，不外推为通用误差规律 & 限定结论 \\
7 & uart\_\allowbreak{}rx 起始位检测涉及首级同步信号 & 同步链健壮性风险 & 没有独立错误率测量；不估计概率或声称风险不存在 & 未关闭 \\
8 & 并行验证共享临时 .so，致 setup FAIL & 测试框架 & 每调用使用独立临时目录；同场景复验通过 & 已修 \\
\bottomrule
\end{longtable}
}

当前服务固件另已修复 UART RX 状态位读取、空输入 ISE padding 与服务帧标志问题，桥 RTL 修复重复制入字节和复位恢复；真实 CPU 联合仿真与后续板测证据分别归档，不仅依赖 host 模型。

\Needspace{85pt}

\paragraph{3.3.3.7 未闭环项}

{\small
\renewcommand{\thetable}{3-12}
\begin{longtable}{@{}>{\raggedright\arraybackslash}p{\dimexpr 0.3168\linewidth-2\tabcolsep\relax}>{\raggedright\arraybackslash}p{\dimexpr 0.3759\linewidth-2\tabcolsep\relax}>{\raggedright\arraybackslash}p{\dimexpr 0.3073\linewidth-2\tabcolsep\relax}@{}}
\caption{验证未闭环项}\label{tab:3-12}\\
\toprule
\textbf{项} & \textbf{状态} & \textbf{影响} \\
\midrule
\endfirsthead
\multicolumn{3}{@{}l}{\footnotesize 续表 3-12}\\
\toprule
\textbf{项} & \textbf{状态} & \textbf{影响} \\
\midrule
\endhead
板上 ISA 抽样回归（含已知 \texttt{rv32ui\_\allowbreak{}\allowbreak{}ma\_\allowbreak{}\allowbreak{}data} 失败） & \blockedflag{未执行} & 不得声称指令集已在板上全覆盖 \\
完整 1,089 组 KAT 板级运行 & 289 组容量内通过；800 组超容量 NOT\_\allowbreak{}RUN & 全量仅 host，当前服务不承诺超容量功能 \\
MMIO 长度能力 & 仅 $\leq$16 B & 不得泛化为任意长度 \\
板级真实 IO 延迟 & \blockedflag{not\_\allowbreak{}\allowbreak{}measured} & 20 ns 为建模上界，见 §3.3.2 \\
PS/PL 与 PYNQ & Notebook、benchmark、9/9 错误门禁已执行；小时级、冷启动重复性未测 & 当前架构无共享可缓存 buffer；不应伪造缓存测试 \\
RV64 / 完整异常矩阵 & 未验收 & 定向中断通过不等于所有 ISE 抢占上下文均通过 \\
芯片丝印速度等级现场核对 & \blockedflag{未做} & 沿用 \texttt{-\allowbreak{}1} 约束，风险记录于 \texttt{docs/\allowbreak{}risks.\allowbreak{}md} \\
\bottomrule
\end{longtable}
}

\evid{证据 docs/\allowbreak{}ascon\_\allowbreak{}standard\_\allowbreak{}baseline.\allowbreak{}md §3-\allowbreak{}-\allowbreak{}§5、docs/\allowbreak{}ascon\_\allowbreak{}key\_\allowbreak{}management.\allowbreak{}md §自动验证入口、project\_\allowbreak{}log/\allowbreak{}actual\_\allowbreak{}progress.\allowbreak{}md 2026-\allowbreak{}10-\allowbreak{}25/\allowbreak{}10-\allowbreak{}26/\allowbreak{}10-\allowbreak{}27 条目、docs/\allowbreak{}freeze\_\allowbreak{}checkpoint\_\allowbreak{}20261027.\allowbreak{}md §3-\allowbreak{}-\allowbreak{}§5}

\Needspace{85pt}

\subsection{性能量化}

\Needspace{85pt}

\paragraph{3.3.4.1 当前 PYNQ 固定输入 warm benchmark}\label{ux5f53ux524d-pynq-ux56faux5b9aux8f93ux5165-warm-benchmark}

输入为锁定官方 KAT \#529，PT=16 B、AD=0；每模式预热 5 次，再各采集 30 次加密和 30 次解密。全部结果与官方 CT、Tag、PT 比对，service simulated=false、selftest\_mask=7。周期测量范围为固件 svc\_aead\_encrypt/decrypt 调用边界，不含 PS 与 UART。

{\footnotesize
\renewcommand{\thetable}{3-13}
\begin{longtable}{@{}>{\raggedright\arraybackslash}p{\dimexpr 0.0931\linewidth-2\tabcolsep\relax}>{\raggedright\arraybackslash}p{\dimexpr 0.1587\linewidth-2\tabcolsep\relax}>{\raggedright\arraybackslash}p{\dimexpr 0.1436\linewidth-2\tabcolsep\relax}>{\raggedright\arraybackslash}p{\dimexpr 0.1512\linewidth-2\tabcolsep\relax}>{\raggedright\arraybackslash}p{\dimexpr 0.1587\linewidth-2\tabcolsep\relax}>{\raggedright\arraybackslash}p{\dimexpr 0.1436\linewidth-2\tabcolsep\relax}>{\raggedright\arraybackslash}p{\dimexpr 0.1512\linewidth-2\tabcolsep\relax}@{}}
\caption{当前 PYNQ 固定输入周期与加速比}\label{tab:3-13}\\
\toprule
\textbf{模式} & \textbf{加密中位 cycles} & \textbf{加密 cycles/B} & \textbf{加密周期加速比} & \textbf{解密中位 cycles} & \textbf{解密 cycles/B} & \textbf{解密周期加速比} \\
\midrule
\endfirsthead
\multicolumn{7}{@{}l}{\footnotesize 续表 3-13}\\
\toprule
\textbf{模式} & \textbf{加密中位 cycles} & \textbf{加密 cycles/B} & \textbf{加密周期加速比} & \textbf{解密中位 cycles} & \textbf{解密 cycles/B} & \textbf{解密周期加速比} \\
\midrule
\endhead
SW & 23342 & 1458.875 & 1.000$\times$ & 24005 & 1500.312 & 1.000$\times$ \\
HW/MMIO & 1439 & 89.938 & 16.221$\times$ & 1388 & 86.750 & 17.295$\times$ \\
ISE & 3995 & 249.688 & 5.843$\times$ & 4385 & 274.062 & 5.474$\times$ \\
\bottomrule
\end{longtable}
}

六档 raw\_cycles 各 30 个，且每档数值完全一致，故 min/median/P95/max 相同。这是固定输入 warm 路径的重复采样结果，不是覆盖全部长度的平均性能，也不是小时级稳定性证据。证据：\texttt{data/\allowbreak{}pynq/\allowbreak{}notebook\_\allowbreak{}20261029\_\allowbreak{}130920/\allowbreak{}pynq\_\allowbreak{}warm\_\allowbreak{}benchmark.\allowbreak{}json}。

\Needspace{85pt}

\paragraph{3.3.4.2 PS 调用墙钟与 CPU 周期的区别}\label{ps-ux8c03ux7528ux5899ux949fux4e0e-cpu-ux5468ux671fux7684ux533aux522b}

同一批采样的 PS 调用墙钟如下，由 raw\_wall\_ns 复算中位数：

{\small
\renewcommand{\thetable}{3-14}
\begin{longtable}{@{}>{\raggedright\arraybackslash}p{\dimexpr 0.2671\linewidth-2\tabcolsep\relax}>{\raggedright\arraybackslash}p{\dimexpr 0.3665\linewidth-2\tabcolsep\relax}>{\raggedright\arraybackslash}p{\dimexpr 0.3665\linewidth-2\tabcolsep\relax}@{}}
\caption{PS 调用中位墙钟}\label{tab:3-14}\\
\toprule
\textbf{模式} & \textbf{加密中位 ms} & \textbf{解密中位 ms} \\
\midrule
\endfirsthead
\multicolumn{3}{@{}l}{\footnotesize 续表 3-14}\\
\toprule
\textbf{模式} & \textbf{加密中位 ms} & \textbf{解密中位 ms} \\
\midrule
\endhead
SW & 5.216 & 4.957 \\
HW/MMIO & 4.479 & 4.262 \\
ISE & 5.154 & 4.986 \\
\bottomrule
\end{longtable}
}

PS 墙钟包括 Python、MMIO 访问、UART 字节传输和轮询等待。CPU 周期下降并未等比例转化为 PS 响应时间下降；不得将 16.221$\times$ 或 5.843$\times$ 标成 Python API 的响应时间加速比。两条测量边界分别保留，不以串口理论传输时间代替实测。

\Needspace{85pt}

\paragraph{3.3.4.3 历史 PL-only 基准的独立对照}\label{ux5386ux53f2-pl-only-ux57faux51c6ux7684ux72ecux7acbux5bf9ux7167}

历史三套基准固件在同一板卡、配置 100 MHz、16 B 加密 warm 条件下经串口输出 mcycle：

{\small
\renewcommand{\thetable}{3-15}
\begin{longtable}{@{}>{\raggedright\arraybackslash}p{\dimexpr 0.2031\linewidth-2\tabcolsep\relax}>{\raggedright\arraybackslash}p{\dimexpr 0.2031\linewidth-2\tabcolsep\relax}>{\raggedright\arraybackslash}p{\dimexpr 0.2485\linewidth-2\tabcolsep\relax}>{\raggedright\arraybackslash}p{\dimexpr 0.3453\linewidth-2\tabcolsep\relax}@{}}
\caption{历史 PL-only 周期对照}\label{tab:3-15}\\
\toprule
\textbf{模式} & \textbf{历史板上 cycles} & \textbf{对应 RTL 仿真 cycles} & \textbf{历史周期加速比} \\
\midrule
\endfirsthead
\multicolumn{4}{@{}l}{\footnotesize 续表 3-15}\\
\toprule
\textbf{模式} & \textbf{历史板上 cycles} & \textbf{对应 RTL 仿真 cycles} & \textbf{历史周期加速比} \\
\midrule
\endhead
SW & 20806 & 17353 & 1.00$\times$ \\
ISE & 2570 & 2496 & 8.10$\times$ \\
MMIO 调用 & 688 & 670 & 30.24$\times$ \\
MMIO 核本体观测 & 103 & 99 & 不与完整 AEAD 调用并列计算加速比 \\
\bottomrule
\end{longtable}
}

历史日志为 \texttt{data/\allowbreak{}ascon/\allowbreak{}results/\allowbreak{}board/\allowbreak{}uart\_\allowbreak{}20261026\_\allowbreak{}sw\_\allowbreak{}and\_\allowbreak{}ise.\allowbreak{}txt} 与 \texttt{uart\_\allowbreak{}20261027\_\allowbreak{}mmio.\allowbreak{}txt}。 688 与 103 的差额为 585 周期，占调用周期约 85.0\%；这仅说明该历史测量中核外开销占比，不直接套用到当前服务固件。板上与仿真差异的完整定量归因未完成，不外推为一般平台规律。当前主表始终使用 §3.3.4.1 的 PYNQ 同次运行结果。

\Needspace{85pt}

\paragraph{3.3.4.4 资源与时序的统计范围}

当前 PYNQ overlay 的资源、WNS、DRC 取自独立 pynq\_wrapper 实现报告（§2.3.3.9）。历史 PL-only 的 12426 LUT、16156 FF、34.5 BRAM tile、4 DSP 为整个 SoC 的实现资源，包含 CPU、Cache、总线、存储与加速逻辑，不是单独密码运算逻辑的面积。 SW/ISE/HW 历史镜像共用同一硬件结构，固件 COE 不同，不能据此推导各方案的边际资源消耗。复用同一个 Verilog 模块定义也不意味着不同实例在运行时共享一块电路。

本文只陈述已有整机实现结果及同平台周期对照，不提供未测的功耗、Fmax 或各加速通路独立面积。

\Needspace{85pt}

\paragraph{3.3.4.5 对比结论的适用边界}

当前结论限于相同 RV32/PYNQ 平台、相同输入与操作、相同固件调用边界的三模式对照。 host 的完整标准验证、历史 PL-only 周期和当前 PYNQ 结果分别成立，不互相替代。超出 16 B 的板级运行不属于当前服务能力，既不计为通过，也不列为本次交付必须完成的新功能。不据整机资源或短消息延迟宣称优于其他器件、其他工具或其他论文实现。

\evid{证据 data/\allowbreak{}pynq/\allowbreak{}notebook\_\allowbreak{}20261029\_\allowbreak{}130920/\allowbreak{}pynq\_\allowbreak{}warm\_\allowbreak{}benchmark.\allowbreak{}json；board/\allowbreak{}vivado\_\allowbreak{}zynq7000\_\allowbreak{}xc7z020\_\allowbreak{}pynq/\allowbreak{}validation\_\allowbreak{}20261029.\allowbreak{}md；data/\allowbreak{}ascon/\allowbreak{}results/\allowbreak{}board/\allowbreak{}uart\_\allowbreak{}20261026\_\allowbreak{}sw\_\allowbreak{}and\_\allowbreak{}ise.\allowbreak{}txt；data/\allowbreak{}ascon/\allowbreak{}results/\allowbreak{}board/\allowbreak{}uart\_\allowbreak{}20261027\_\allowbreak{}mmio.\allowbreak{}txt}

\Needspace{85pt}

\subsection{完成度与判据矩阵}

{\small
\renewcommand{\thetable}{3-16}
\begin{longtable}{@{}>{\raggedright\arraybackslash}p{\dimexpr 0.2343\linewidth-2\tabcolsep\relax}>{\raggedright\arraybackslash}p{\dimexpr 0.3662\linewidth-2\tabcolsep\relax}>{\raggedright\arraybackslash}p{\dimexpr 0.3995\linewidth-2\tabcolsep\relax}@{}}
\caption{完成度与判据矩阵}\label{tab:3-16}\\
\toprule
\textbf{工作项} & \textbf{当前状态} & \textbf{证据与边界} \\
\midrule
\endfirsthead
\multicolumn{3}{@{}l}{\footnotesize 续表 3-16}\\
\toprule
\textbf{工作项} & \textbf{当前状态} & \textbf{证据与边界} \\
\midrule
\endhead
RV32 平台与三方案运行 & 已有实板证据 & 当前 pynq-svc-v1 在同一 overlay 下切换 SW/HW/ISE \\
Final 标准一致性（host） & 1089 组通过 & 完整 host 结果不等于完整板级结果 \\
当前板级容量内 KAT & 三模式各 289/289 & PT/AD 各 0--16 B，错误 Tag 各 289/289 \\
当前 warm benchmark & 已执行 & KAT \#529、预热 5 次、每档 30 次，原始 cycles 与墙钟齐全 \\
Notebook 从头顺序执行 & 已有运行结果 & warm benchmark 已执行本体尚未归档；不能写成未执行 \\
板级错误恢复 & 9/9 checks PASS & 7/7 reset 恢复，每次三模式 KAT 通过 \\
连续请求 & 已验证有限次数 & 每模式 50 次加密与 50 次解密；非小时级验证 \\
当前实现时序 / DRC & 已有报告 & WNS +0.034 ns、TNS 0、DRC 0 Error/0 Critical；Warning 非零 \\
独立 RO-PUF/DNA 原型 & 已实现并有一次板级自检通过 & v3 quality=128/128，同次重建状态通过；非 PYNQ release，见扩展原型章节 \\
完整 1089 组板测 & NOT\_\allowbreak{}RUN & 800 组超接口容量；本稿不承诺全量板级支持 \\
板上 ISA 抽样 & NOT\_\allowbreak{}RUN & 历史仿真 49/50，ma\_\allowbreak{}data 因非对齐访存失败；不称全套 ISA 通过 \\
RV64、完整异常与抢占矩阵 & 未验收 & 不作为当前 RV32 bare-metal 演示的已验证能力 \\
小时级资源泄漏与冷启动重复性 & not\_\allowbreak{}measured & 无 RSS/FD 长时曲线，不以有限次数请求替代 \\
独立 FCLK 与完整实物取证 & 部分缺失 & 已有正面照片；芯片完整丝印、跳线与 SD 镜像等仍待归档 \\
实际交付 Skill & rtl-\allowbreak{}design-\allowbreak{}discipline 与 pynq-\allowbreak{}overlay-\allowbreak{}workflow 已成形 & 后者工具层 47 项测试、13/13 变异检出通过；智能体对照与实板未运行，不称效果已验证（见 §2.5.3） \\
clean-clone / 最终包 & 未完成 & 已能使用 release 不等于已在干净环境端到端重建 \\
\bottomrule
\end{longtable}
}

\Needspace{85pt}

\section{复现说明}

\Needspace{85pt}

\subsection{环境与版本}

以下步骤从空白环境开始，按''主机验证 $\rightarrow$ 仿真回归 $\rightarrow$ 构建比特流 $\rightarrow$ 上板''顺序执行。每一步都有可判定的通过标准，失败时保留日志并停止，不跳到下一步。所有命令在仓库根目录执行。

{\small
\renewcommand{\thetable}{3-17}
\begin{longtable}{@{}>{\raggedright\arraybackslash}p{\dimexpr 0.2220\linewidth-2\tabcolsep\relax}>{\raggedright\arraybackslash}p{\dimexpr 0.7780\linewidth-2\tabcolsep\relax}@{}}
\caption{环境与版本}\label{tab:3-17}\\
\toprule
\textbf{项} & \textbf{版本或要求} \\
\midrule
\endfirsthead
\multicolumn{2}{@{}l}{\footnotesize 续表 3-17}\\
\toprule
\textbf{项} & \textbf{版本或要求} \\
\midrule
\endhead
源码 & \texttt{git clone https:/\allowbreak{}/\allowbreak{}github.\allowbreak{}com/\allowbreak{}HChenZ/\allowbreak{}amd\_\allowbreak{}\allowbreak{}fpga\_\allowbreak{}\allowbreak{}2026.\allowbreak{}git}，记录 \texttt{git rev-\allowbreak{}parse HEAD} 与 \texttt{git status -\allowbreak{}-\allowbreak{}short} \\
EDA & Vivado 2025.2（\texttt{build.\allowbreak{}tcl} 会检查版本，不符即退出） \\
固件工具链 & RV32 GCC（归档构建为 13.2.0，\texttt{rv32im\_\allowbreak{}\allowbreak{}zicsr} / \texttt{ilp32}） \\
主机与仿真 & Python 3、Icarus Verilog；主机标准验证可在 WSL/Linux 运行 \\
板卡 & 正点原子领航者 ZYNQ7020，PYNQ 2.7.0 镜像，网口与电脑同网段 \\
\bottomrule
\end{longtable}
}

\Needspace{180pt}

\subsection{主机验证与仿真回归}

\begin{codeblock}[samepage=true]

python3 sim/ascon/run_validation.py
python3 sim/ascon/check_failure_contract.py

# 2. RTL 回归：独立向量 + 生产源集门禁
python3 sim/ascon/prepare_rtl_vectors.py
python sim/ascon/run_rtl_regression.py

# 3. SoC 回归：真实 CPU、三模式固件、mailbox 与 PYNQ 桥
python -B sim/board/run_tests.py --mode regression --pynq-p0
python sim/pynq/run_bridge.py
python sim/pynq/run_service.py
\end{codeblock}

通过标准为各命令 exit 0，且生成的 JSON 汇总中无 FAIL。归档结果见 \texttt{build/\allowbreak{}reports/\allowbreak{}pynq\_\allowbreak{}final\_\allowbreak{}20261028/\allowbreak{}}。

\Needspace{128pt}

\subsection{固件与比特流构建}

\begin{codeblock}[samepage=true]
# 1. 编译服务固件
#    产物：build/local/pynq_svc/pynq_svc.{elf,bin,coe}
bash src/firmware/pynq_svc/build.sh

# 2. 构建独立 PYNQ overlay
#    不打开、不修改历史冻结工程
vivado -mode batch -source board/vivado_zynq7000_xc7z020_pynq/build.tcl
\end{codeblock}

固件反汇编中应能找到 custom-0 指令（P12、P8），证明 ISE 确实进入二进制。构建完成后核对时序报告 WNS $\geq$ 0、TNS = 0、DRC 无 Error。复建的 bitstream 不要求与发布版逐字节相同，应分别核对源输入、功能、时序、资源与新产物哈希；直接复现已发布结果时使用 \texttt{release/\allowbreak{}} 中的配套文件。

\Needspace{206pt}

\subsection{上板运行}

写卡、串口跳线与启动模式按 \texttt{board/\allowbreak{}vivado\_\allowbreak{}zynq7000\_\allowbreak{}xc7z020\_\allowbreak{}pynq/\allowbreak{}handoff\_\allowbreak{}20261028.\allowbreak{}md} 第 C 节执行。登录板卡后：

\begin{codeblock}[samepage=true]
cd ~/jupyter_notebooks
git clone https://github.com/HChenZ/amd_fpga_2026.git
cd amd_fpga_2026
# bit/hwh/固件必须配套
sha256sum -c board/vivado_zynq7000_xc7z020_pynq/release/sha256sums.txt
# 三模式各 289 组 KAT 与坏 Tag 拒绝
python3 src/pynq/board_demo.py --full-kat --out pynq_board_full.json
# KAT #529，预热 5 次，采样 30 次
python3 src/pynq/board_benchmark.py --out pynq_warm_benchmark.json
# 9 项错误恢复，每模式 50 次往返
python3 src/pynq/board_error_gate.py --out pynq_board_error_gate.json
\end{codeblock}

交互演示使用 \texttt{src/\allowbreak{}pynq/\allowbreak{}pynq\_\allowbreak{}board\_\allowbreak{}demo.\allowbreak{}ipynb}（修改首个单元的仓库路径后 Restart Kernel $\rightarrow$ Run All）； CLI 与 Notebook 共享同一通道，不要同时运行。通过标准：boot 自检 mask=7、固件标识 pynq-svc-v1、 simulated=false，三模式 KAT 与坏 Tag 拒绝全部通过；周期结果应与 1.4 节主要性能指标同量级。

\Needspace{85pt}

\subsection{复现边界}

\begin{itemize}
\tightlist
\item
  干净环境下从克隆到上板的完整端到端复现尚未由第三方执行；本节命令均为本项目实际使用过的入口。
\item
  PT/AD 超过 16 B 的 800 组官方 KAT 超出当前接口容量，板上不执行；完整 1089 组只在主机验证。
\item
  当前时序余量 34 ps，换用其他 Vivado 版本、IP 版本或约束需重新签核。
\item
  板卡速度等级与独立频率未实测；独立 RO-PUF 原型不在 PYNQ 发布版中，其复现见 §2.3.4 所列工程。
\end{itemize}
